\par\bigskip
\hypertarget{native-384-ceres2}{}
\begin{center}
\pdfbookmark[2]{2. The Perturbations of the Planets}{ceres-002}[2.]\par
\medskip
\textls[120]{The Perturbations of the Planets.}
\end{center}
\[
\begin{gathered}
\left.\begin{array}{c}x,y,z\\x',y',z'\end{array}\right\}
\text{ coordinates of the }
\left\{\begin{array}{l}\text{perturbed planet}\\\text{perturbing planet}\end{array}\right\},\\[4pt]
\mu=\frac{\text{mass of the perturbing pl.}}{\text{mass of the Sun}}
\end{gathered}
\]
\[
\begin{gathered}
r=\BareRoot{(xx+yy+zz)},\qquad
r'r'=\BareRoot{(x'x'+y'y'+z'z')},\\
\mathfrak{R}=\text{dist[ance] of the two planets [from each other]}
\end{gathered}
\]
\noindent\(m,A\) mean [longitude] and [mean] distance of any plan[et from the Sun]
\[
\begin{aligned}
R&=\frac{xx'+yy'+zz'}{r'^{3}}-\frac{1}{\mathfrak{R}}\\
 &=\frac{r}{r'r'}\cos(\text{ang. dist. of planets})-\frac{1}{\mathfrak{R}}\mathpunct{\cdot}
\end{aligned}
\]
\SourceLeaf{385}
\par\medskip
\noindent\hspace*{\parindent}\textls[120]{Fundamental equations:}
\[
\begin{aligned}
0&=\frac{\mathrm{d}\mathrm{d}x}{\mathrm{d}m^{2}}+\frac{xA^{3}}{r^{3}}+\mu A^{3}\left(\frac{\mathrm{d}R}{\mathrm{d}x}\right)\\[4pt]
0&=\frac{\mathrm{d}\mathrm{d}y}{\mathrm{d}m^{2}}+\frac{yA^{3}}{r^{3}}+\mu A^{3}\left(\frac{\mathrm{d}R}{\mathrm{d}y}\right)\\[4pt]
0&=\frac{\mathrm{d}\mathrm{d}z}{\mathrm{d}m^{2}}+\frac{zA^{3}}{r^{3}}+\mu A^{3}\left(\frac{\mathrm{d}R}{\mathrm{d}z}\right)
\end{aligned}
\]
\[
\begin{aligned}
[1]\quad 0&=\frac{\mathrm{d}r^{2}+rr\mathrm{d}v^{2}}{\mathrm{d}m^{2}}-\frac{2A^{3}}{r}+\frac{A^{3}}{a}+2\mu A^{3}\int\partial R\\[6pt]
[2]\quad 0&=\frac{\mathrm{d}r^{2}+r\mathrm{d}\mathrm{d}r}{\mathrm{d}m^{2}}-\frac{\mathrm{d}r^{2}+rr\mathrm{d}v^{2}}{\mathrm{d}m^{2}}\\
&\qquad+\frac{A^{3}}{r}+\mu A^{3}r\left(\frac{\mathrm{d}R}{\mathrm{d}r}\right)\\[6pt]
[3]\quad 0&=\frac{r\mathrm{d}\mathrm{d}r+\mathrm{d}r^{2}}{\mathrm{d}m^{2}}-\frac{A^{3}}{r}+\frac{A^{3}}{a}\\
&\qquad+2\mu A^{3}\int\partial R+\mu A^{3}r\left(\frac{\mathrm{d}R}{\mathrm{d}r}\right)\mathpunct{\cdot}
\end{aligned}
\]
\par
\noindent\hspace*{\parindent}We set
\[
2\mu A^{3}\int\partial R+\mu A^{3}r\left(\frac{\mathrm{d}R}{\mathrm{d}r}\right)=Q;
\]
[the expression \(\partial R\) is to be understood as varying only the coordinates of the perturbed planet in the differentiation.
\par
\noindent\hspace*{\parindent}If \(\delta r\) denotes the perturbations of \(r\), then equation 3 gives, on neglecting the 2nd powers of the perturbing mass]
\[
[4]\qquad 0=\frac{\mathrm{d}\mathrm{d}(r\delta r)}{\mathrm{d}m^{2}}+\frac{A^{3}}{r^{3}}r\delta r+Q.
\]
\par
\noindent\hspace*{\parindent}[To introduce the true longitude \(v\) as the independent variable,] the full equation is thus written, with \(\delta r=\rho\),
\[
\begin{aligned}
[5]\quad 0&=\frac{r\mathrm{d}\mathrm{d}\rho+2\mathrm{d}r\mathrm{d}\rho+\rho\mathrm{d}\mathrm{d}r}{\mathrm{d}m^{2}}\\
&\qquad-(r\mathrm{d}\rho+\rho\mathrm{d}r)\frac{\mathrm{d}\mathrm{d}m}{\mathrm{d}m^{3}}+\frac{A^{3}}{r^{3}}r\rho+Q.
\end{aligned}
\]
\par
\noindent\hspace*{\parindent}Now, setting the eccentric[ity] of the perturbed pl.\(=e\) and \(\BareRoot{(1-ee)}=\gamma\), we have [to the same degree of accuracy]
\[
rr\mathrm{d}v=\gamma A^{\frac32}a^{\frac12}\mathrm{d}m
\]
and, treating \(\mathrm{d}v\) as constant,
\[
2r\mathrm{d}r\mathrm{d}v=\gamma A^{\frac32}a^{\frac12}\mathrm{d}\mathrm{d}m
\]
[so equation 5 becomes]
\[
[6]\qquad 0=\frac{r\mathrm{d}\mathrm{d}\rho+\rho\mathrm{d}\mathrm{d}r-\frac{2\rho}{r}\mathrm{d}r^{2}}{\mathrm{d}m^{2}}+\frac{A^{3}}{r^{3}}r\rho+Q.
\]
\par
\noindent\hspace*{\parindent}Finally, we have [again to the same degree of accuracy]
\SourceLeaf{386}
\[
\begin{gathered}
\frac{\mathrm{d}\mathrm{d}\dfrac1r}{\mathrm{d}v^{2}}+\frac1r=\frac1{a\gamma\gamma}\qquad\text{or}\\[4pt]
\frac{\mathrm{d}v^{2}}{r}-\frac{\mathrm{d}\mathrm{d}r}{rr}+\frac{2\mathrm{d}r^{2}}{r^{3}}=\frac{\mathrm{d}v^{2}}{a\gamma\gamma}\mathpunct{\cdot}
\end{gathered}
\]
\par
\noindent\hspace*{\parindent}Thus
\[
[7]\qquad 0=\frac{r\mathrm{d}\mathrm{d}\rho+r\rho\mathrm{d}v^{2}-\frac{rr\rho}{a\gamma\gamma}\mathrm{d}v^{2}}{\mathrm{d}m^{2}}+\frac{A^{3}}{r^{3}}r\rho+Q
\]
or [since the third and fourth terms on the right cancel]
\[
0=\frac{\mathrm{d}\mathrm{d}\rho}{\mathrm{d}v^{2}}+\rho+\frac{\mathrm{d}m^{2}}{\mathrm{d}v^{2}}\frac{Q}{r}
\]
or
\[
[8]\qquad 0=\frac{\mathrm{d}\mathrm{d}\rho}{\mathrm{d}v^{2}}+\rho+\frac{r^{3}Q}{A^{3}a\gamma\gamma}\mathpunct{\cdot}
\]
\par
\noindent\hspace*{\parindent}It follows at once that
\[
\rho=\cos v\int\frac{r^{3}Q}{A^{3}a\gamma\gamma}\sin v\,\mathrm{d}v-\sin v\int\frac{r^{3}Q}{A^{3}a\gamma\gamma}\cos v\,\mathrm{d}v
\]
or [since \(rr\mathrm{d}v=\gamma aa\mathrm{d}(nt)\)]
\[
*\quad\rho=\left\{\begin{aligned}&\cos v\int rQ\sin v\,\mathrm{d}(nt)\\[-2pt]&-\sin v\int rQ\cos v\,\mathrm{d}(nt)\end{aligned}\right\}\frac{a}{A^{3}\BareRoot{(1-ee)}}\mathpunct{\cdot}
\]
\par\medskip
\begin{center}\rule{.27\linewidth}{.4pt}\end{center}
\par
\noindent\hspace*{\parindent}[Further, it follows from equations 1 and 3 that]
\[
\begin{aligned}
0&=\frac{\frac32\mathrm{d}r^{2}-\frac12rr\mathrm{d}v^{2}+2r\mathrm{d}\mathrm{d}r}{\mathrm{d}m^{2}}-\frac{A^{3}}{r}+\frac32\frac{A^{3}}{a}\\
&\qquad+3\mu A^{3}\int\partial R+2\mu A^{3}r\left(\frac{\mathrm{d}R}{\mathrm{d}r}\right)
\end{aligned}
\]
[thus, after introducing \(\rho\) and \(\delta v\), cancelling the unperturbed values, and neglecting the 2nd powers of the perturbations,]
\[
\begin{aligned}
0&=\frac{3\mathrm{d}r\mathrm{d}\rho-rr\mathrm{d}v\mathrm{d}(\delta v)-r\rho\mathrm{d}v^{2}+2\rho\mathrm{d}\mathrm{d}r+2r\mathrm{d}\mathrm{d}\rho}{\mathrm{d}m^{2}}\\
&\qquad+\frac{A^{3}\rho}{rr}+3\mu A^{3}\int\partial R+2\mu A^{3}r\left(\frac{\mathrm{d}R}{\mathrm{d}r}\right)\\[6pt]
[9]\quad&=\frac{\mathrm{d}(\rho\mathrm{d}r+2r\mathrm{d}\rho)-rr\mathrm{d}v\mathrm{d}(\delta v)-r\rho\mathrm{d}v^{2}+\rho\mathrm{d}\mathrm{d}r}{\mathrm{d}m^{2}}\\
&\qquad+\frac{A^{3}\rho}{rr}+3\mu A^{3}\int\partial R+2\mu A^{3}r\left(\frac{\mathrm{d}R}{\mathrm{d}r}\right)\mathpunct{\cdot}
\end{aligned}
\]
\par
\noindent\hspace*{\parindent}Now, by the above equation [2, to sufficient accuracy,]
\[
0=\frac{-r\mathrm{d}v^{2}+\mathrm{d}\mathrm{d}r}{\mathrm{d}m^{2}}+\frac{A^{3}}{rr}\mathpunct{\text{\raisebox{.45ex}{,}}}
\]
so
\[
\begin{aligned}
[10]\quad 0&=\frac{\mathrm{d}(\rho\mathrm{d}r+2r\mathrm{d}\rho)-rr\mathrm{d}v\mathrm{d}(\delta v)}{\mathrm{d}m^{2}}\\
&\qquad+3\mu A^{3}\int\partial R+2\mu A^{3}r\left(\frac{\mathrm{d}R}{\mathrm{d}r}\right)\mathpunct{\cdot}
\end{aligned}
\]
\par
\noindent\hspace*{\parindent}Consequently
\[
\begin{aligned}
*\quad\delta v&=\frac{\rho\mathrm{d}r+2r\mathrm{d}\rho}{aa\BareRoot{(1-ee)}\mathord{.}\mathrm{d}(nt)}\\
&\quad+\frac{3\mu a\int\!\int\partial R\mathord{.}\mathrm{d}(nt)+2\mu a\int r\left(\frac{\mathrm{d}R}{\mathrm{d}r}\right)\mathrm{d}(nt)}{\BareRoot{(1-ee)}}\mathpunct{\cdot}
\end{aligned}
\]
\par\medskip
\begin{center}\rule{.27\linewidth}{.4pt}\end{center}
\SourceLeaf{387}
\par
\noindent\hspace*{\parindent}Finally, if we take the original orbit of the [perturbed] planet as the fundamental plane in which the coordinates \(x,y\) lie and set \(z=r\theta\), then
\[
0=\frac{\mathrm{d}\mathrm{d}(r\theta)}{\mathrm{d}m^{2}}+\frac{A^{3}\theta}{rr}+\mu A^{3}\left(\frac{\mathrm{d}R}{\mathrm{d}z}\right)\mathpunct{\text{\raisebox{.45ex}{,}}}
\]
which [equation] can be integrated in the same way as [4] above, hence
\[
*\quad\theta=\left\{\begin{aligned}&\cos v\int r\left(\frac{\mathrm{d}R}{\mathrm{d}z}\right)\sin v\mathord{.}\mathrm{d}(nt)\\[2pt]&-\sin v\int r\left(\frac{\mathrm{d}R}{\mathrm{d}z}\right)\cos v\mathord{.}\mathrm{d}(nt)\end{aligned}\right\}\frac{\mu a}{\BareRoot{(1-ee)}}\mathpunct{\cdot}
\]
\par
\noindent\hspace*{\parindent}[\(\theta\) denotes the latitude, or the sine of the latitude, above the unperturbed orbital plane, taken as fixed.] Thus
\[
\begin{gathered}
\text{the hel[iocentric] latitude is increased by }\\[2pt]
\frac{\cos(\mathrm{incl.})}{\cos(\mathrm{hel.\ lat.})}\theta\\[6pt]
\text{the hel[iocentric] longitude is decreased by }\\[2pt]
\frac{\theta\sin(\mathrm{incl.})\cos(\mathrm{arg.\ lat.})}{\cos(\mathrm{hel.\ lat.})^{2}}\mathpunct{\cdot}
\end{gathered}
\]
\par\medskip
\begin{center}\rule{.27\linewidth}{.4pt}\end{center}
